% =============================================================================
% 2. BACKGROUND AND RELATED WORK
% Target: 2 pages (~1000 words)
% Source: tese/cap2_reduced.tex (compressed) + gap synthesis from SLR (cap. 3)
% Status: SKELETON
% =============================================================================

\section{Background and Related Work}
\label{sec:background}

% ---------- §2.1 Process mining fundamentals ----------
\subsection{Process Mining}
\label{sec:bg-pm}

\textbf{[TODO §2.1 — 0.5 p.]}
Process mining~\cite{aalst2016process} extracts process knowledge from
event logs. The discipline distinguishes three families of techniques:
discovery (e.g., the Inductive Miner~\cite{leemans2013discovering}),
conformance checking (alignments), and enhancement.
The Inductive Miner is adopted here for its
formal soundness guarantee on the discovered process tree, which is
required by Stage~2 of the pipeline (Section~\ref{sec:stages-s2}).

% Key concepts to summarize:
% - event log (case id, activity, timestamp)
% - process model (Petri net / process tree)
% - fitness, precision, generalization (four-quality framework)

% ---------- §2.2 Markov chains for business processes ----------
\subsection{Absorbing Markov Chains}
\label{sec:bg-markov}

\textbf{[TODO §2.2 — 0.5 p.]}
Absorbing Markov chains provide closed-form expressions for expected
absorption times and absorption probabilities via the fundamental matrix
$N = (I-Q)^{-1}$~\cite{kemeny1976finite}. Their application to business
process analysis has been explored by Rogge-Solti and
Weske~\cite{RoggeSolti2015NonMarkovianSPN} via stochastic Petri nets,
and by~\cite{Spedicato2017} via direct estimation from event logs.
PATHCAST adopts the latter approach but couples it with empirical
sojourn-time distributions, avoiding the parametric assumptions of
classical semi-Markov treatments.

% ---------- §2.3 Stochastic forecasting in software engineering ----------
\subsection{Stochastic Forecasting in Software Engineering}
\label{sec:bg-forecasting}

\textbf{[TODO §2.3 — 0.5 p.]}
Throughput-based Monte Carlo, popularized by
Vacanti~\cite{vacanti2015} and extended
by~\cite{savage2009flaw, magennis2015forecasting}, samples future
completion dates from the empirical throughput distribution of a team.
The approach is operationally simple and widely adopted in agile
practice, but treats the process as a black box: rework, state-specific
delays, and path diversity are not captured.
Recent work on probabilistic process forecasting in business processes
includes~\cite{polato2018time, verenich2019survey}; none of these
methods is specialised to the SDLC nor integrates the heterogeneous
data sources (Git, issue tracker, CI/CD, code review) characteristic of
software repositories.

% ---------- §2.4 Gap synthesis ----------
\subsection{Pipeline-Fragmentation Gap}
\label{sec:bg-gap}

\textbf{[TODO §2.4 — 0.5 p.]}
Table~\ref{tab:gap-synthesis} summarizes the coverage of process mining,
stochastic modeling, and forecasting across five representative
categories of work in the SDLC literature. Only a single primary study
in the SLR of~\cite{almeida2026slr} satisfies all three dimensions,
motivating the integrated pipeline proposed in this paper.

\begin{table}[ht]
\centering
\caption{Coverage of PM, stochastic modeling, and forecasting across
representative SDLC-related research families. ``\checkmark'' indicates
explicit treatment; ``--'' indicates absence.}
\label{tab:gap-synthesis}
\small
\begin{tabular}{@{}lccc l@{}}
\toprule
\textbf{Family} & \textbf{PM} & \textbf{Stochastic} &
\textbf{Forecasting} & \textbf{Representative} \\
\midrule
SDLC process discovery       & \checkmark & --          & --          & \cite{rubin2007process} \\
Throughput Monte Carlo       & --         & \checkmark  & \checkmark  & \cite{vacanti2015}     \\
Effort/lead-time ML          & --         & --          & \checkmark  & \cite{choetkiertikul2018deep} \\
Stochastic Petri nets (BPM)  & \checkmark & \checkmark  & --          & \cite{RoggeSolti2015NonMarkovianSPN} \\
\textbf{PATHCAST (this paper)} & \checkmark & \checkmark & \checkmark & --- \\
\bottomrule
\end{tabular}
\end{table}
